% ৪.x HTML List
\section{HTML List}

\begin{learningbox}
এই টপিক শেষে শিক্ষার্থী পারবে:
\begin{itemize}
\item HTML list কী তা ব্যাখ্যা করতে।
\item Ordered List, Unordered List ও Description List-এর পার্থক্য লিখতে।
\item \texttt{<ol>}, \texttt{<ul>}, \texttt{<li>}, \texttt{<dl>}, \texttt{<dt>} ও \texttt{<dd>} tag ব্যবহার করতে।
\item \texttt{type} ও \texttt{start} attribute-এর কাজ বুঝতে।
\item nested list তৈরি করতে।
\item HSC practical ও board exam-এর list-related প্রশ্নের উত্তর দিতে।
\end{itemize}
\end{learningbox}

\subsection{HTML List কী?}

\begin{definitionbox}{HTML List}
HTML List হলো web page-এ একাধিক তথ্য বা item-কে ধারাবাহিক, bullet বা description আকারে সাজিয়ে প্রদর্শনের পদ্ধতি।
\end{definitionbox}

Subject list, food menu, book list, routine, product feature, instruction step, website menu বা glossary - এসব তথ্য list আকারে দেখালে user সহজে পড়তে ও বুঝতে পারে। HTML-এ list ব্যবহার করলে content organized হয় এবং web page দেখতে পরিষ্কার লাগে।

\begin{boardbox}{পরীক্ষার সংক্ষিপ্ত উত্তর}
HTML list হলো web page-এ একাধিক item সংগঠিতভাবে প্রদর্শনের পদ্ধতি। Ordered List, Unordered List ও Description List - HTML-এর প্রধান list type।
\end{boardbox}


\subsection{HTML List-এর প্রকারভেদ}

\begin{center}
\begin{tabular}{|p{0.25\textwidth}|p{0.18\textwidth}|p{0.43\textwidth}|}
\hline
\textbf{List Type} & \textbf{Tag} & \textbf{ব্যবহার} \\
\hline
Ordered List & \texttt{<ol>} & সংখ্যা, অক্ষর বা roman marker সহ ক্রমানুসার list \\
\hline
Unordered List & \texttt{<ul>} & bullet/symbol সহ list; item-এর order গুরুত্বপূর্ণ নয় \\
\hline
Description List & \texttt{<dl>} & term ও description আকারে list \\
\hline
\end{tabular}
\end{center}

\begin{keypointbox}{মনে রাখো}
\texttt{<ol>} ও \texttt{<ul>} list-এর প্রতিটি item লেখার জন্য \texttt{<li>} tag ব্যবহার করা হয়। কিন্তু \texttt{<dl>} list-এ \texttt{<li>} নয়; term-এর জন্য \texttt{<dt>} এবং description-এর জন্য \texttt{<dd>} ব্যবহার করা হয়।
\end{keypointbox}

\subsection{Ordered List: \texttt{<ol>}}

\begin{definitionbox}{Ordered List}
Ordered List হলো এমন list যেখানে item গুলো নির্দিষ্ট ক্রমে সংখ্যা, অক্ষর বা roman number দিয়ে প্রদর্শিত হয়।
\end{definitionbox}

যখন list-এর item গুলোর order বা ক্রম গুরুত্বপূর্ণ, তখন ordered list ব্যবহার করা হয়। যেমন: instruction step, exam routine, ranking, chapter list, admission process বা algorithm steps।

\subsubsection{Syntax}

\begin{examplebox}{Code}
\begin{verbatim}
<ol>
  <li>First item</li>
  <li>Second item</li>
  <li>Third item</li>
</ol>
\end{verbatim}
\end{examplebox}

\subsubsection{Example}

\begin{examplebox}{Code}
\begin{verbatim}
<!DOCTYPE html>
<html>
<body>
  <h2>HSC ICT Chapter List</h2>

  <ol>
    <li>ICT: World and Bangladesh Perspective</li>
    <li>Communication System and Networking</li>
    <li>Number System and Digital Device</li>
    <li>Web Design Introduction and HTML</li>
  </ol>
</body>
</html>
\end{verbatim}
\end{examplebox}

\begin{examplebox}{Output ও ব্যাখ্যা}
\textbf{Output:}
\begin{enumerate}
\item ICT: World and Bangladesh Perspective
\item Communication System and Networking
\item Number System and Digital Device
\item Web Design Introduction and HTML
\end{enumerate}
\textbf{ব্যাখ্যা:} \texttt{<ol>} ordered list শুরু করে এবং প্রতিটি \texttt{<li>} একটি করে list item তৈরি করে। Browser defaultভাবে marker হিসেবে 1, 2, 3 ব্যবহার করে।
\end{examplebox}

\subsection{Ordered List-এর Attribute}

\subsubsection{\texttt{type} attribute}

\texttt{type} attribute দিয়ে ordered list-এর marker কীভাবে দেখাবে তা নির্ধারণ করা যায়।

\begin{center}
\begin{tabular}{|p{0.25\textwidth}|p{0.35\textwidth}|}
\hline
\textbf{Type} & \textbf{Output marker} \\
\hline
\texttt{type="1"} & 1, 2, 3 \\
\hline
\texttt{type="A"} & A, B, C \\
\hline
\texttt{type="a"} & a, b, c \\
\hline
\texttt{type="I"} & I, II, III \\
\hline
\texttt{type="i"} & i, ii, iii \\
\hline
\end{tabular}
\end{center}

\begin{examplebox}{Code}
\begin{verbatim}
<ol type="A">
  <li>HTML</li>
  <li>CSS</li>
  <li>JavaScript</li>
</ol>
\end{verbatim}
\end{examplebox}

\begin{examplebox}{Output ও ব্যাখ্যা}
\textbf{Output:}\\
A. HTML\\
B. CSS\\
C. JavaScript\\
\textbf{ব্যাখ্যা:} \texttt{type="A"} দেওয়ায় marker number না হয়ে capital letter হিসেবে দেখায়।
\end{examplebox}

\subsubsection{\texttt{start} attribute}

\texttt{start} attribute দিয়ে ordered list কোন number থেকে শুরু হবে তা নির্ধারণ করা যায়।

\begin{examplebox}{Code}
\begin{verbatim}
<ol start="5">
  <li>HTML</li>
  <li>CSS</li>
  <li>JavaScript</li>
</ol>
\end{verbatim}
\end{examplebox}

\begin{examplebox}{Output ও ব্যাখ্যা}
\textbf{Output:}
\begin{enumerate}[start=5]
\item HTML
\item CSS
\item JavaScript
\end{enumerate}
\textbf{ব্যাখ্যা:} \texttt{start="5"} দেওয়ায় list marker ১ থেকে না শুরু হয়ে ৫ থেকে শুরু হয়েছে।
\end{examplebox}

\begin{keypointbox}{Spec note}
\texttt{reversed} attribute দিলে ordered list উল্টো ক্রমে গণনা হয়।
আর কোনো নির্দিষ্ট item-এর numbering আলাদা করতে সেই \texttt{<li>} element-এ \texttt{value} attribute ব্যবহার করা যায়।
\end{keypointbox}

\begin{boardbox}{Ordered List: exam answer}
Ordered List হলো এমন HTML list যেখানে item গুলো number, letter বা roman numeral দিয়ে ক্রমানুসারে প্রদর্শিত হয়। Ordered list তৈরির জন্য \texttt{<ol>} tag এবং প্রতিটি item-এর জন্য \texttt{<li>} tag ব্যবহার করা হয়।
\end{boardbox}

\subsection{Unordered List: \texttt{<ul>}}

\begin{definitionbox}{Unordered List}
Unordered List হলো এমন list যেখানে item গুলো bullet বা symbol দিয়ে প্রদর্শিত হয় এবং item-এর নির্দিষ্ট ক্রম গুরুত্বপূর্ণ নয়।
\end{definitionbox}

যখন item-এর order গুরুত্বপূর্ণ নয়, তখন unordered list ব্যবহার করা হয়। যেমন: menu list, feature list, subject list, shopping list, notice points, important topic list বা website navigation menu।

\subsubsection{Syntax}

\begin{examplebox}{Code}
\begin{verbatim}
<ul>
  <li>Item one</li>
  <li>Item two</li>
  <li>Item three</li>
</ul>
\end{verbatim}
\end{examplebox}

\subsubsection{Example}

\begin{examplebox}{Code}
\begin{verbatim}
<!DOCTYPE html>
<html>
<body>
  <h2>Web Design Tools</h2>

  <ul>
    <li>HTML</li>
    <li>CSS</li>
    <li>JavaScript</li>
    <li>Browser</li>
  </ul>
</body>
</html>
\end{verbatim}
\end{examplebox}

\begin{examplebox}{Output ও ব্যাখ্যা}
\textbf{Output:}
\begin{itemize}
\item HTML
\item CSS
\item JavaScript
\item Browser
\end{itemize}
\textbf{ব্যাখ্যা:} \texttt{<ul>} unordered list তৈরি করে। Browser defaultভাবে bullet marker দেখায়।
\end{examplebox}

\subsection{Unordered List-এর \texttt{type} Attribute}

HSC book-style HTML code-এ unordered list-এর marker পরিবর্তনের জন্য \texttt{type} attribute দেখা যায়।

\begin{center}
\begin{tabular}{|p{0.30\textwidth}|p{0.34\textwidth}|}
\hline
\textbf{Type} & \textbf{Output idea} \\
\hline
\texttt{type="disc"} & filled bullet \\
\hline
\texttt{type="circle"} & circle bullet \\
\hline
\texttt{type="square"} & square bullet \\
\hline
\end{tabular}
\end{center}

\begin{examplebox}{Code}
\begin{verbatim}
<ul type="square">
  <li>Monitor</li>
  <li>Keyboard</li>
  <li>Mouse</li>
</ul>
\end{verbatim}
\end{examplebox}

\begin{examplebox}{Output ও ব্যাখ্যা}
\textbf{Output:}\\
\(\blacksquare\) Monitor\\
\(\blacksquare\) Keyboard\\
\(\blacksquare\) Mouse\\
\textbf{ব্যাখ্যা:} \texttt{type="square"} দেওয়ায় list marker square bullet আকারে দেখা যাবে।
\end{examplebox}

\begin{mistakebox}{Modern note}
Modern HTML/CSS-এ list marker design করার জন্য CSS ব্যবহার করা বেশি standard। তবে HSC exam-এর জন্য \texttt{type="disc"}, \texttt{type="circle"} ও \texttt{type="square"} জানা দরকার।
\end{mistakebox}

\begin{boardbox}{Unordered List: exam answer}
Unordered List হলো এমন HTML list যেখানে item গুলো bullet symbol দিয়ে প্রদর্শিত হয়। এটি তৈরির জন্য \texttt{<ul>} tag এবং প্রতিটি item-এর জন্য \texttt{<li>} tag ব্যবহার করা হয়।
\end{boardbox}

\subsection{List Item: \texttt{<li>}}

\begin{definitionbox}{List item}
\texttt{<li>} tag list-এর প্রতিটি item লিখতে ব্যবহৃত হয়।
\end{definitionbox}

\texttt{<ol>} এবং \texttt{<ul>} - দুই ধরনের list-এর item লিখতেই \texttt{<li>} tag ব্যবহার করা হয়।

\begin{examplebox}{Code}
\begin{verbatim}
<ul>
  <li>Bangla</li>
  <li>English</li>
  <li>ICT</li>
</ul>
\end{verbatim}
\end{examplebox}

\begin{examplebox}{Output ও ব্যাখ্যা}
\textbf{Output:}
\begin{itemize}
\item Bangla
\item English
\item ICT
\end{itemize}
\textbf{ব্যাখ্যা:} প্রতিটি \texttt{<li>} tag একটি আলাদা list item তৈরি করেছে।
\end{examplebox}

\begin{keypointbox}{HSC note}
কিছু পুরনো example-এ \texttt{<li>} closing tag না দিলেও চলে দেখানো হতে পারে। কিন্তু বই লেখার সময় ও practical coding-এ \texttt{<li>Item</li>} আকারে closing tag দিয়ে লিখলে code পরিষ্কার ও standard হয়।
\end{keypointbox}

\subsection{Description List: \texttt{<dl>}, \texttt{<dt>}, \texttt{<dd>}}

\begin{definitionbox}{Description List}
Description List হলো এমন list যেখানে কোনো term বা বিষয়ের নাম এবং তার description বা ব্যাখ্যা একসাথে দেখানো হয়।
\end{definitionbox}

Glossary, definition list, FAQ, product specification, contact details, dictionary-like content বা term and meaning দেখাতে description list ব্যবহার করা যায়।

\begin{center}
\begin{tabular}{|p{0.18\textwidth}|p{0.58\textwidth}|}
\hline
\textbf{Tag} & \textbf{কাজ} \\
\hline
\texttt{<dl>} & Description list শুরু ও শেষ করে \\
\hline
\texttt{<dt>} & Term বা name নির্দেশ করে \\
\hline
\texttt{<dd>} & Term-এর description নির্দেশ করে \\
\hline
\end{tabular}
\end{center}

\subsubsection{Syntax}

\begin{examplebox}{Code}
\begin{verbatim}
<dl>
  <dt>Term</dt>
  <dd>Description of the term</dd>
</dl>
\end{verbatim}
\end{examplebox}

\subsubsection{Example}

\begin{examplebox}{Code}
\begin{verbatim}
<!DOCTYPE html>
<html>
<body>
  <h2>Web Terms</h2>

  <dl>
    <dt>HTML</dt>
    <dd>HTML is used to create the structure of a web page.</dd>

    <dt>CSS</dt>
    <dd>CSS is used to design and style a web page.</dd>

    <dt>JavaScript</dt>
    <dd>JavaScript is used to add interactivity to a web page.</dd>
  </dl>
</body>
</html>
\end{verbatim}
\end{examplebox}

\begin{examplebox}{Output ও ব্যাখ্যা}
\textbf{Output:}\\
\textbf{HTML}\\
\hspace*{1cm}HTML is used to create the structure of a web page.\\
\textbf{CSS}\\
\hspace*{1cm}CSS is used to design and style a web page.\\
\textbf{JavaScript}\\
\hspace*{1cm}JavaScript is used to add interactivity to a web page.\\
\textbf{ব্যাখ্যা:} \texttt{<dt>} tag term দেখায় এবং \texttt{<dd>} tag সেই term-এর description দেখায়।
\end{examplebox}

\begin{boardbox}{Description List: exam answer}
Description List হলো এমন HTML list যেখানে term ও তার description দেখানো হয়। এটি তৈরির জন্য \texttt{<dl>} tag, term লেখার জন্য \texttt{<dt>} tag এবং description লেখার জন্য \texttt{<dd>} tag ব্যবহার করা হয়।
\end{boardbox}

\begin{keypointbox}{Spec note}
একটি \texttt{dl} element-এ একাধিক \texttt{dt} দিয়ে একটি name-এর একাধিক term দেওয়া যেতে পারে এবং একাধিক \texttt{dd} দিয়ে
একটি term-এর একাধিক description দেওয়া যেতে পারে। প্রয়োজন হলে পুরো group-কে \texttt{div} দিয়ে wrap করা যায়,
বিশেষ করে styling বা metadata-এর জন্য।
\end{keypointbox}

\subsection{Nested List}

\begin{definitionbox}{Nested List}
একটি list-এর item-এর ভিতরে আরেকটি list ব্যবহার করলে তাকে Nested List বলা যায়।
\end{definitionbox}

\begin{examplebox}{Code}
\begin{verbatim}
<ul>
  <li>Science
    <ol>
      <li>Physics</li>
      <li>Chemistry</li>
      <li>Biology</li>
    </ol>
  </li>

  <li>Business Studies
    <ol>
      <li>Accounting</li>
      <li>Finance</li>
      <li>Management</li>
    </ol>
  </li>
</ul>
\end{verbatim}
\end{examplebox}

\begin{examplebox}{Output ও ব্যাখ্যা}
\textbf{Output:}
\begin{itemize}
\item Science
  \begin{enumerate}
  \item Physics
  \item Chemistry
  \item Biology
  \end{enumerate}
\item Business Studies
  \begin{enumerate}
  \item Accounting
  \item Finance
  \item Management
  \end{enumerate}
\end{itemize}
\textbf{ব্যাখ্যা:} বাইরের listটি \texttt{<ul>} এবং প্রতিটি group-এর ভিতরে \texttt{<ol>} list আছে। এইভাবে category ও subcategory দেখানো যায়।
\end{examplebox}

\begin{keypointbox}{Nested list কোথায় ব্যবহার হয়}
Subject group, website menu, chapter and sub-topic, category and subcategory, course module বা product feature group দেখাতে nested list ব্যবহার করা যায়।
\end{keypointbox}

\subsection{Ordered, Unordered ও Description List-এর পার্থক্য}

\begin{center}
\small
\begin{tabular}{|p{0.18\textwidth}|p{0.22\textwidth}|p{0.22\textwidth}|p{0.22\textwidth}|}
\hline
\textbf{বিষয়} & \textbf{Ordered List} & \textbf{Unordered List} & \textbf{Description List} \\
\hline
Tag & \texttt{<ol>} & \texttt{<ul>} & \texttt{<dl>} \\
\hline
Item tag & \texttt{<li>} & \texttt{<li>} & \texttt{<dt>}, \texttt{<dd>} \\
\hline
Marker & সংখ্যা/অক্ষর/roman & bullet/symbol & term-description \\
\hline
Order গুরুত্বপূর্ণ? & হ্যাঁ & না & term অনুযায়ী \\
\hline
ব্যবহার & step, ranking, routine & topic, feature, menu & definition, glossary \\
\hline
Example & 1, 2, 3 & bullet & HTML - markup language \\
\hline
\end{tabular}
\normalsize
\end{center}

\subsection{Complete Example Code}

\begin{examplebox}{Code}
\begin{verbatim}
<!DOCTYPE html>
<html>
<head>
  <meta charset="UTF-8">
  <title>HTML List Example</title>
</head>
<body>

  <h1>HTML List Example</h1>

  <h2>Ordered List</h2>
  <ol type="1">
    <li>Open Notepad</li>
    <li>Write HTML code</li>
    <li>Save as index.html</li>
    <li>Open with browser</li>
  </ol>

  <h2>Unordered List</h2>
  <ul type="square">
    <li>HTML</li>
    <li>CSS</li>
    <li>JavaScript</li>
  </ul>

  <h2>Description List</h2>
  <dl>
    <dt>HTML</dt>
    <dd>HTML creates the structure of a web page.</dd>

    <dt>CSS</dt>
    <dd>CSS designs the appearance of a web page.</dd>
  </dl>

</body>
</html>
\end{verbatim}
\end{examplebox}

\begin{examplebox}{Output ধারণা}
Browser-এ প্রথমে ordered list 1, 2, 3, 4 আকারে দেখা যাবে। এরপর unordered list square bullet আকারে দেখা যাবে। শেষে description list term ও description আকারে দেখা যাবে।
\end{examplebox}

\begin{mistakebox}{সাধারণ ভুল}
\begin{itemize}
\item \texttt{<ol>} ও \texttt{<ul>} গুলিয়ে ফেলা।
\item \texttt{<li>} tag না দেওয়া।
\item \texttt{<dl>} tag-এর ভিতরে \texttt{<li>} ব্যবহার করা।
\item \texttt{<dt>} ও \texttt{<dd>} ভুলভাবে ব্যবহার করা।
\item ordered list-এ \texttt{type} value ভুল লেখা।
\item \texttt{<ul>} tag-এ number marker আশা করা।
\item list-এর closing tag না দেওয়া।
\item nested list ভুলভাবে close করা।
\end{itemize}
\end{mistakebox}

\begin{boardbox}{Correct nested structure}
\begin{examplebox}{Code}
\begin{verbatim}
<ul>
  <li>ICT
    <ol>
      <li>HTML</li>
      <li>Networking</li>
    </ol>
  </li>
</ul>
\end{verbatim}
\end{examplebox}
\end{boardbox}

\subsection{পরীক্ষার জন্য সংক্ষিপ্ত উত্তর}

HTML list ব্যবহার করে web page-এর তথ্য সুন্দর ও সংগঠিতভাবে প্রদর্শন করা যায়। HTML-এ প্রধানত তিন ধরনের list ব্যবহৃত হয় - Ordered List, Unordered List এবং Description List। Ordered List-এ item গুলো সংখ্যা, অক্ষর বা roman number দিয়ে প্রদর্শিত হয় এবং এটি \texttt{<ol>} tag দিয়ে তৈরি করা হয়। Unordered List-এ item গুলো bullet symbol দিয়ে প্রদর্শিত হয় এবং এটি \texttt{<ul>} tag দিয়ে তৈরি করা হয়। Ordered ও unordered list-এর প্রতিটি item লেখার জন্য \texttt{<li>} tag ব্যবহার করা হয়। Description List-এ term ও description দেখানোর জন্য \texttt{<dl>}, \texttt{<dt>} এবং \texttt{<dd>} tag ব্যবহার করা হয়।

\subsection{ছোট প্রশ্নোত্তর}

\begin{enumerate}
\item HTML list কী?\\
\textbf{উত্তর:} Web page-এ একাধিক item সংগঠিতভাবে প্রদর্শনের পদ্ধতিকে HTML list বলে।

\item HTML-এ list কয় প্রকার?\\
\textbf{উত্তর:} সাধারণত ৩ প্রকার - Ordered List, Unordered List, Description List।

\item Ordered List তৈরির tag কোনটি?\\
\textbf{উত্তর:} \texttt{<ol>}।

\item Unordered List তৈরির tag কোনটি?\\
\textbf{উত্তর:} \texttt{<ul>}।

\item List item লেখার tag কোনটি?\\
\textbf{উত্তর:} \texttt{<li>}।

\item Description List তৈরির tag কোনটি?\\
\textbf{উত্তর:} \texttt{<dl>}।

\item Description List-এ term লেখার tag কোনটি?\\
\textbf{উত্তর:} \texttt{<dt>}।

\item Description List-এ description লেখার tag কোনটি?\\
\textbf{উত্তর:} \texttt{<dd>}।

\item \texttt{<ol type="A">} দিলে output কেমন হবে?\\
\textbf{উত্তর:} A, B, C আকারে list marker দেখাবে।

\item \texttt{<ul type="square">} দিলে output কেমন হবে?\\
\textbf{উত্তর:} square bullet আকারে item দেখাবে।
\end{enumerate}

\subsection{Exam Focus MCQ}

\begin{enumerate}
\item HTML-এ list সাধারণত কয় প্রকার?\\
ক) ২ \quad খ) ৩ \quad গ) ৪ \quad ঘ) ৫

\item Ordered List তৈরির tag কোনটি?\\
ক) \texttt{<ul>} \quad খ) \texttt{<ol>} \quad গ) \texttt{<li>} \quad ঘ) \texttt{<dl>}

\item Unordered List তৈরির tag কোনটি?\\
ক) \texttt{<ul>} \quad খ) \texttt{<ol>} \quad গ) \texttt{<dt>} \quad ঘ) \texttt{<dd>}

\item List item যুক্ত করার tag কোনটি?\\
ক) \texttt{<ol>} \quad খ) \texttt{<li>} \quad গ) \texttt{<ul>} \quad ঘ) \texttt{<dl>}

\item Description List-এ term লেখার tag কোনটি?\\
ক) \texttt{<dd>} \quad খ) \texttt{<dt>} \quad গ) \texttt{<li>} \quad ঘ) \texttt{<ul>}

\item \texttt{<ol type="i">} দিলে list marker কেমন হবে?\\
ক) A, B, C \quad খ) a, b, c \quad গ) I, II, III \quad ঘ) i, ii, iii

\item Unordered list-এর marker square করতে কোনটি ব্যবহার করা যায়?\\
ক) \texttt{<ul type="square">} \quad খ) \texttt{<ol type="square">} \quad গ) \texttt{<li type="A">} \quad ঘ) \texttt{<dl type="square">}
\end{enumerate}

\begin{boardbox}{Answer key}
১-খ, ২-খ, ৩-ক, ৪-খ, ৫-খ, ৬-ঘ, ৭-ক
\end{boardbox}
